% ===========================================================================
\chapter{Setting up}\label{ch:setup}

Make a folder for the project and go into it. Every file in this tutorial
goes in this one folder.

\begin{shell}
mkdir datalog
cd datalog
\end{shell}

Copy the four sample programs, \code{social.dl}, \code{flights.dl},
\code{access.dl} and \code{chain.dl}, from the \code{examples/datalog}
folder of the Nex distribution into it. They are plain text; here is the
first:

\begin{shell}
% social.dl: who follows whom, and what follows from it.

follows(ann, bo).
follows(bo, ann).
follows(bo, cy).
follows(cy, dee).
follows(dee, eli).
follows(ann, fay).
follows(fay, cy).

% Two people are friends when each follows the other.
friends(X, Y) :- follows(X, Y), follows(Y, X).

% Whose posts can reach X, if everyone reshares what they see: everyone X
% follows, everyone they follow, and so on.
sees(X, Y) :- follows(X, Y).
sees(X, Z) :- follows(X, Y), sees(Y, Z).

% People who follow someone in common.
in_common(X, Y) :- follows(X, Z), follows(Y, Z), X != Y.

% Someone X might like to follow: followed by someone X follows, but not
% yet by X, and not X.
suggest(X, Z) :- follows(X, Y), follows(Y, Z), X != Z, not follows(X, Z).

?- sees(ann, Who).
?- suggest(ann, Who).
\end{shell}

A \texttt{\%} starts a comment that runs to the end of the line. Names that
start with a lower-case letter, like \code{ann} and \code{follows}, are
constants and predicates; names that start with a capital letter, like
\code{X} and \code{Who}, are variables. \code{X != Y} is a test: the rule
for \code{in_common} only applies when \code{X} and \code{Y} are different
people, or everyone would have something in common with themselves.
\code{not follows(X, Z)} holds when there is no such fact, so a suggestion
is only made to someone who doesn't already follow that person.

\section{The files}

The engine is split into one file per idea, each building on the ones
before it with \code{intern}, which makes everything one Nex file defines
available in another:

\begin{center}
\begin{tabular}{@{}ll@{}}
\toprule
\code{terms.nex} & terms, atoms, literals, rules and programs \\
\code{lexer.nex} & cutting text into tokens \\
\code{parser.nex} & building a program from the tokens \\
\code{bindings.nex} & giving variables values; matching \\
\code{database.nex} & storing facts \\
\code{evaluate.nex} & deriving every fact the rules imply \\
\code{query.nex} & answering questions \\
\code{analysis.nex} & checking a program; negation in layers \\
\code{explain.nex} & saying why a fact holds \\
\code{datalog.nex} & the command line \\
\code{checks.nex} & checks of all of the above \\
\bottomrule
\end{tabular}
\end{center}

Along the way we write small throwaway programs, called \code{try_}
something, to see each piece work. Run any of them with \code{nex}:

\begin{shell}
nex try_terms.nex
\end{shell}

A word about printing. Nex's \code{print} shows a string with its quotes,
which is handy when debugging. Every Nex program also has a ready-made
\code{console}, and \code{console.print_line("hello")} prints the text as
it is, on a line of its own. The try programs use whichever suits
them.

\section{Contracts}

Nex checks every \code{require}, \code{ensure} and invariant as the program
runs. The engine is full of them, and several do real work: they are how
we will know that the fast way of evaluating a program finds exactly what
the slow way finds. Some are expensive. When you want speed rather than
checking, run a program with \code{--skip-contracts}:

\begin{shell}
nex datalog.nex chain.dl --skip-contracts
\end{shell}
